\chapter*{Introduction Générale}
\addcontentsline{toc}{chapter}{Introduction Générale}

L'intelligence artificielle générative a connu une évolution rapide : les modèles de langage de grande taille, initialement utilisés comme moteurs de génération ou d'assistance conversationnelle, sont désormais intégrés à des systèmes capables de \emph{raisonner sur une tâche}, de sélectionner des outils, d'interroger des ressources externes et d'exécuter des séquences d'actions. Cette évolution donne naissance aux \textbf{agents IA autonomes}. Un agent peut, par exemple, lire un document, rechercher une information dans une base, appeler une API, créer un ticket, envoyer un message ou manipuler un environnement de développement. L'intérêt opérationnel est considérable, mais cette autonomie modifie profondément le modèle de menace traditionnel.

Dans une application classique, le développeur contrôle directement la chaîne qui mène d'une entrée à une action. Dans un système agentique, une partie de cette chaîne dépend du raisonnement produit par le LLM. Le modèle n'exécute généralement pas lui-même la fonction : il propose un appel structuré, puis l'application déclenche l'outil correspondant. OpenRouter décrit ce mécanisme de \textit{tool calling} comme une boucle dans laquelle le modèle suggère un appel d'outil, l'application exécute la fonction, puis retourne son résultat au modèle afin qu'il poursuive sa tâche \cite{openroutertools}. Cette séparation est fondamentale pour la sécurité, car elle crée une frontière à laquelle une décision d'autorisation peut être imposée.

Le problème apparaît lorsque le modèle traite simultanément des \textbf{instructions légitimes} et du \textbf{contenu non fiable}. Un utilisateur peut demander de résumer un document, tandis que ce document contient lui-même une instruction cachée demandant à l'agent de lire une ressource confidentielle et de l'envoyer vers une destination externe. Ce phénomène est connu sous le nom de \textit{prompt injection indirecte}. Les travaux de Greshake et al. ont montré que le mélange entre données et instructions dans des applications intégrant des LLM ouvre des scénarios de détournement à distance \cite{greshake2023}. Des benchmarks plus récents, comme InjecAgent et AgentDojo, confirment que les attaques sont réelles mais fortement dépendantes du modèle, de la tâche, de l'attaque et du contexte d'exécution \cite{injecagent,agentdojo}.

Le \textbf{Model Context Protocol (MCP)} ajoute une nouvelle dimension à ce problème. MCP est un standard ouvert destiné à connecter les applications d'IA aux systèmes dans lesquels résident les données et les outils. Les serveurs MCP peuvent exposer des outils, des ressources et d'autres capacités à des clients agentiques. La révision 2026-07-28 du protocole poursuit notamment l'évolution vers un cœur sans état, le routage par en-têtes, un modèle d'extensions et un renforcement de l'autorisation \cite{mcp2026}. Cette standardisation facilite l'interopérabilité, mais elle agrandit également la surface de confiance : l'agent doit désormais raisonner sur des capacités qui peuvent provenir de plusieurs serveurs, propriétaires ou niveaux de sensibilité.

Dans ce contexte, notre projet \Aegis{} cherche à répondre à la question suivante :

\begin{quote}
\textbf{Comment permettre à un agent IA autonome d'utiliser des outils externes via MCP, tout en empêchant qu'une instruction non fiable, un outil malveillant, un accès excessif ou un raisonnement erroné ne conduise à l'exécution d'une action non autorisée ?}
\end{quote}

La réponse proposée repose sur une idée volontairement simple : \textbf{le LLM propose, Aegis décide}. Le modèle reste responsable du raisonnement fonctionnel et peut demander des actions. En revanche, l'autorisation finale est externalisée vers un mécanisme déterministe nommé \Guard{}. Celui-ci évalue notamment l'identité de l'agent, le but initial de l'utilisateur, la provenance de l'instruction, le niveau de confiance de la source, la sensibilité des données, le risque de l'outil, la destination d'une éventuelle sortie d'information et la présence ou non d'une autorisation explicite. Une décision peut alors être prise : autoriser, refuser, réduire la portée, mettre en quarantaine ou demander une validation humaine.

Le projet comporte également deux autres dimensions. La première est offensive : \Red{} fournit un cadre de Red Teaming destiné à reproduire de manière contrôlée des attaques comme l'injection de prompt directe ou indirecte, l'empoisonnement du RAG, la corruption de mémoire, le \textit{tool poisoning}, l'abus d'identité ou l'exfiltration par chaîne d'outils. La seconde est organisationnelle : \GRC{} transforme les résultats techniques en risques, contrôles, preuves, indicateurs et décisions de traitement. Cette dimension s'appuie notamment sur le NIST AI Risk Management Framework, ISO/IEC 42001, ISO/IEC 23894 et certains principes d'ISO/IEC 27001 \cite{nistairmf,iso42001,iso23894}. Le contexte marocain est pris en compte à travers la loi 09-08 et les recommandations de la CNDP relatives à la protection des données à caractère personnel \cite{cndp0908,cndpconditions}.

Le présent rapport suit la structure recommandée dans le cadre pédagogique du projet. Le \textbf{chapitre 1} introduit le contexte général, l'étude de l'existant, la problématique, les objectifs ainsi que la méthodologie de travail. Le \textbf{chapitre 2} présente l'état de l'art des agents IA, du MCP, des principales menaces agentiques et des référentiels de sécurité et de gouvernance. Le \textbf{chapitre 3} détaille la conception fonctionnelle et sécuritaire d'AegisMCP, les besoins, les modèles UML, le modèle de menace et la logique de décision. Le \textbf{chapitre 4} décrit l'infrastructure et l'architecture proposée, ainsi que les choix technologiques. Enfin, le \textbf{chapitre 5} présente la réalisation du prototype, les modules effectivement développés à la date de rédaction, les premiers scénarios expérimentaux, les résultats observés et les étapes nécessaires pour atteindre l'évaluation finale.

Une règle méthodologique guide l'ensemble du rapport : \textbf{aucun résultat de sécurité n'est inventé}. Les résultats quantitatifs présentés correspondent uniquement aux expériences réellement exécutées. Lorsqu'une fonctionnalité appartient à l'architecture cible mais n'est pas encore implémentée à la date de cette version, elle est explicitement décrite comme telle. Cette distinction permet de conserver la valeur scientifique du travail et d'éviter de confondre conception, hypothèse et résultat mesuré.
